\documentclass[titlepage,12pt]{article}

\usepackage[margin=0.5in,headsep=2mm]{geometry}
\usepackage{mathptmx}
\usepackage{csquotes}
\usepackage{setspace}
\usepackage{fancyhdr}
\usepackage[style=chem-acs]{biblatex}
\usepackage{subcaption}
\usepackage{tikz}
\usepackage{chemfig}
\usepackage{mhchem}
\usepackage[hidelinks]{hyperref}

\MakeOuterQuote{"}

\pagestyle{fancy}
\fancyhf{}
\fancyhead[R]{\thepage}
\renewcommand{\headrulewidth}{0pt}

\addbibresource{../main.bib}
\DefineBibliographyStrings{english}{bibliography={References}}

\usetikzlibrary{decorations.pathmorphing}
\colorlet{blx}{blue!40!black!80}

\setchemfig{atom sep=2em,fixed length=true,bond offset=3pt,cram width=3pt}

\title{A Review of Schrock's Abiological Nitrogen Fixation Cycle}
\author{
    Steven Labalme\\
    Department of Chemistry, University of Chicago\\
    CHEM 20200: Inorganic Chemistry II\\
    Dr. John Anderson
}
\date{26 May 2021}

\begin{document}




\begin{doublespace}
    \maketitle
\end{doublespace}



Nitrogen must be fixed in order to be biologically useful, and nature does not fix enough nitrogen on its own to sustain the current human population. Thus, in order to fertilize humanity’s food, the Haber-Bosch process is employed to artificially fix nitrogen. However, since it is only one (albeit extremely effective) process, there is great interest in understanding other potential methods of nitrogen fixation. The first big step came when researchers synthesized some of the first abiotic transition metal dinitrogen complexes, loosely modeling them after nitrogenase enzymes\textsuperscript{[}\autocite{bib:AllenSenoff}\textsuperscript{]\,[}\autocite{bib:Yamamoto}\textsuperscript{]}. From there, potential catalytic intermediates were characterized, but a complete catalytic reaction remained elusive. However, Yandulov and Schrock outline in \emph{Catalytic Reduction of Dinitrogen to Ammonia at a Single Molybdenum Center} their discovery\textsuperscript{[}\autocite{bib:SchrockYandulov}\textsuperscript{]} of the first catalytic system with observable intermediates analogous to those proposed by Chatt\textsuperscript{[}\autocite{bib:Chatt}\textsuperscript{]}.

\begin{figure}[h!]
    \centering
    \begin{subfigure}[b]{0.49\linewidth}
        \centering
        \chemfig{-[,,,,{decorate,decoration={snake,segment length=2.1mm}}]-[4,0.4,,,opacity=0]\phantom{i}-[2]**6(--(-?[1]-[::-30,0.8](<[::-100,0.6]iPr)-[::30,0.8]-[::30,0.8](-[::-30]iPr)-[::120,0.8]-[::30,0.8]?[1](<:[::-50,0.6]iPr))--(-?[2]-[::-30,0.8](<:[::-100,0.6]iPr)-[::30,0.8]-[::30,0.8](-[::-30]iPr)-[::120,0.8]-[::30,0.8]?[2](<[::-50,0.6]iPr))--)}
        \begin{tikzpicture}[remember picture,overlay]
            \draw [rotate=30] (-0.25,3.78) ellipse (6mm and 2.5mm);
            \draw [rotate=-30] (-7.3,-0.6) ellipse (6mm and 2.5mm);
        \end{tikzpicture}
        \caption{HIPT ligand.}
        \label{fig:catalysta}
    \end{subfigure}
    \begin{subfigure}[b]{0.49\linewidth}
        \centering
        \chemfig{Mo?[1](-N(-[:53]HIPT)-[::-74]-[::-63]-[::-63]N?[1,,{stealth-}]?[2])(<:[:160]N(-[::-53]HIPT)-[::74]-[::67,1.3]?[2])(<[:-150]N(-[::-75,,,,{white,ultra thick,double=black,double distance=3pt}]HIPT)-[::30,,,,{white,ultra thick,double=black,double distance=3pt}]@{C1}-[::91,,,,line width=3pt]?[2,5]@{C2})-[2,,,,{stealth-}]N~[2]N}
        \chemmove{
            \fill
                (C1) circle (1.5pt)
                ([xshift=0.17pt,yshift=0pt]C2.center) circle (1.59pt)
            ;
        }
        \caption{The full monomeric species.}
        \label{fig:catalystb}
    \end{subfigure}
    \caption{The catalyst.}
    \label{fig:catalyst}
\end{figure}

Having studied nitrogen fixation for 20 years, Yandulov and Schrock suspected that the key to a well-defined molecular nitrogen fixation catalyst lay in high oxidation state molybdenum and tungsten species. In order to isolate and study some of these high oxidation state intermediates, the previously semi-successful triamidoamine ligand \ce{[(ArNCH2CH2)3N]^3-} (see Figure \ref{fig:catalystb}) was modified to incorporate the extremely bulky ligand motif HIPT (hexa-iso-propyl-terphenyl or \ce{{3,5-(2,4,6-}iPr3C6H2)2C6H3}; see Figure \ref{fig:catalysta}). These ligands, they hoped, would sterically inhibit formation of previously observed, unreactive bimetallic \ce{N2}-bridging species. Indeed, they were able to synthesize from \ce{\textbf{Mo}Cl} (where $\ce{\textbf{Mo}}=\ce{[HIPTN3N]Mo}$) the six boxed intermediates in Figure \ref{fig:mechanism}, as characterized by NMR and x-ray studies\textsuperscript{[}\autocite{bib:Intermediates146}\textsuperscript{]\,[}\autocite{bib:Intermediates235}\textsuperscript{]}. To support these "intermediates'" relevance to the predicted Chatt-type mechanism, they were additionally able to show that some of them are interconvertible (i.e., could be prepared from others).\par
% As for optimizing the initial results, using a stronger reducing agent (decamethyl chromocene instead of cobaltocene) enabled the reduction of \ce{[\textbf{Mo}(NH3)]+}, a step that must be completed before returning to the \ce{\textbf{Mo}(N2)} resting state. However, to prevent the stronger reducing agent from interacting more with the proton source than the catalyst intermediates, it was also necessary to choose heptane as a solvent (a substance in which the proton source is sparingly soluble but the \ce{BAr^F{}_4} salt intermediates were more soluble, thus, to reiterate, ensuring that the reagents interacted with the catalytic intermediates instead of each other).\par
With this result, Yandulov and Schrock introduced the first well-defined abiological nitrogen-fixation system with relatively high yields and observable, interconvertible intermediates. Additionally, to eliminate the possibility that their nitrogen-containing tetradentate \ce{[HIPTN3N]^3-} ligand or acid was furnishing the nitrogen that was moving around the intermediates, they conducted a study with radiolabeled dinitrogen. Having established that all intermediates and products contained only \ce{{}^15N} at the proposed active site pocket, they concluded that no isotopically labeled nitrogen was exchanged between the HIPT ligand and acid, and the substrate, as would be expected in a true catalytic mechanism. Thus, they reported that their \ce{[HIPTN3N]Mo} species was (1) accepting nitrogen from the atmosphere, (2) likely forming well-defined, isolable intermediates in a plausible theoretical catalytic cycle, and (3) converting between said intermediates.

\begin{figure}[h!]
    \centering
    \setcharge{extra sep=2.5pt}
    \begin{tikzpicture}[remember picture,overlay,scale=1.16,xshift=-1mm]
        \draw [blx,thick] (0.4,-0.2) rectangle ++(3.2,0.9);
        \draw [blx,thick] (5,-0.2) rectangle ++(4.05,0.75);
        \draw [blx,thick] (10.6,-1.1) rectangle ++(4.8,2.4);
        \draw [blx,thick] (11.9,-12.3) rectangle ++(2.4,0.75);
        \draw [blx,thick] (11,-15.2) rectangle ++(3.95,1.1);
        \draw [blx,thick] (0,-8.4) rectangle ++(4.1,2.6);
    \end{tikzpicture}
    \schemestart
        \chemfig{\mathbf{Mo}^{III}-\charge{90[yshift=1mm]=$\oplus$}{N}~\charge{0=\:}{N}}
        \arrow{->[\small\ce{H+, e-}]}
        \chemfig{\mathbf{Mo}^{IV}-\charge{90=\:}{N}=\charge{90=\:}{N}-H}
        \arrow{->[\small\ce{H+}]}
        \chemleft{[}
            \chemfig{\mathbf{Mo}^{VI}=\charge{90=\:}{N}-\charge{90=\:}{N}(-[1]H)(-[7]H)}
        \chemright{]^+}
        \arrow{->[*{0}\small\ce{e-}]}[-90,1.2]
        \chemfig{\mathbf{Mo}^{V}=\charge{90=\:}{N}-\charge{90=\:}{N}(-[1]H)(-[7]H)}
        \arrow{->[*{0}\small\ce{H+}]}[-90,1.3]
        \chemleft{[}
            \chemfig{\mathbf{Mo}^{V}=\charge{90=\:}{N}-\charge{135[yshift=1mm]=$\oplus$}{N}(-[:60]H)(<[:-60]H)(<:[:-20]H)}
        \chemright{]^+}
        \arrow{->[*{0}\small\ce{e-}][*{0}\small\ce{-NH3}]}[-90,1.3]
        \chemfig{\mathbf{Mo}^{VI}~\charge{0=\:}{N}}
        \arrow{->[*{0}\small\ce{H+}]}[-90,1.2]
        \chemleft{[}
            \chemfig{\mathbf{Mo}^{VI}~\charge{90[yshift=1mm]=$\oplus$}{N}-H}
        \chemright{]^+}
        \arrow{->[*{0.-90}\small\ce{e-}]}[180,1.3]
        \chemfig{\mathbf{Mo}^{V}=\charge{90=\:}{N}-H}
        \arrow{->[*{0.-90}\small\ce{H+}]}[180,1.3]
        \chemleft{[}
            \chemfig{\mathbf{Mo}^{V}=\charge{90[yshift=1mm]=$\oplus$}{N}(-[1]H)(-[7]H)}
        \chemright{]^+}
        \arrow{->[*{0}\small\ce{e-}]}[90]
        \chemfig{\mathbf{Mo}^{IV}-\charge{90=\:}{N}(-[1]H)(-[7]H)}
        \arrow{->[*{0}\small\ce{H+}]}[90]
        \chemleft{[}
            \chemfig{\mathbf{Mo}^{IV}-\charge{135[yshift=1mm]=$\oplus$}{N}(-[:60]H)(<[:-60]H)(<:[:-20]H)}
        \chemright{]^+}
        \arrow{->[*{0}\small\ce{e-}]}[90]
        \chemfig{\mathbf{Mo}^{III}-\charge{135[yshift=1mm]=$\oplus$}{N}(-[:60]H)(<[:-60]H)(<:[:-20]H)}
        \arrow{->[*{0}\small\ce{N2}][*{0}\small\ce{-NH3}]}[90]
    \schemestop
    \caption{The mechanism. All boxed intermediates were isolated and characterized. The remaining intermediates are proposed to exist by Chatt, but have yet to be detected.}
    \label{fig:mechanism}
\end{figure}

In introducing the first well-defined catalytic system for nitrogen fixation with observable Chatt-type intermediates, Yandulov and Schrock demonstrated that catalytic nitrogen fixation likely is possible, thus motivating continued research in this field. Although their present system does not approach the turnover numbers and frequencies necessary for it to be industrially viable, it represents a validation of Chatt's preexisting theory and exhibits the potential of high oxidation state transition metals as well as sterically bulky ligands for nitrogen fixation purposes. Additionally, if other, more recent results are anything to go by, the problems of rate and efficiency are likely solvable: For example, Anderson et al. originally reported seven turnovers for a similar nitrogen fixation system, but follow up studies that further purified their reagents and also tinkered with the catalyst and conditions brought the system up to nearly 100 turnovers\textsuperscript{[}\autocite{bib:AndersonPeters}\textsuperscript{]}. However, a somewhat problematic hole in Yandulov and Schrock's results is their inability to characterize Chatt's other five proposed intermediates (see the unboxed species in Figure \ref{fig:mechanism}), especially considering the fact that they were able to characterize species as unstable and rare as the diazine adduct. Their verification of the interconvertibility of their characterized intermediates does support the existence of the other proposed intermediates, but it does not rule out the possibility of something else going on.


% Toward the end of their paper, Yandulov and Schrock noted that the question remains as to whether or not their system is limited to four turnovers. If other, more recent results are anything to go by, it probably is not (for example, Anderson et al. originally reported seven turnovers for a similar nitrogen fixation system, but follow up studies that further purified their reagents and also tinkered with the catalyst and conditions to bring the system up to nearly 100 turnovers\textsuperscript{[}\autocite{bib:AndersonPeters}\textsuperscript{]}). As such, there is likely ample room for improvement. Follow up studies should focus on the conversion from \ce{\textbf{Mo}(NH3)} to \ce{\textbf{Mo}(N2)}, as that was the most difficult step of the catalytic cycle, likely due to the inferiority of dinitrogen as a ligand. However, perhaps the catalyst could be modified in such a way that an \ce{NH3} ligand forces a shape change that makes \ce{NH3} binding much less favorable. Conformation changes are critical to the functioning of most enzymes, so it stands to reason that, if properly utilized, they could also aid in abiological systems.
\newpage



\printbibliography




\end{document}